\section{Discovery by connected cancellation profiles}
\label{sec:cancellation-clusters}

The cancellation parameter can be refined by retaining the positions of the
exceptional edges.  The refinement below is independent of any proposed
global bound on that parameter.  It does not assume that a quadrilateral
sector is supplied, and it accommodates repeated local occurrences of a
global edge.

Let $\mathcal T$ have $t$ tetrahedra and $m$ global edges, one global vertex,
and first Betti number one.  Assume that its underlying compact connected
orientable manifold has exactly one torus boundary.  Write $c$ for a
primitive integral cocycle and $b_e=|c_e|$.  Every essential normal disc has
weights $w=b+u$, where all $u_e$ are nonnegative even integers.  Its
cancellation support is $F=\{e:u_e>0\}$.

Define $G$ to have the global edges as its vertices, with two distinct
vertices adjacent when their edges occur in a common tetrahedron.  Let
\[
 d=\max_e\#\{\Delta:e\text{ occurs in }\Delta\},\qquad
 q(F)=\#\pi_0(G[F]).
\]
An edge occurring several times in one tetrahedron contributes one to
$d$, although all its local occurrences are retained in the equations.
Thus the maximum graph degree is at most $5d$.  Put $q(\varnothing)=0$.
Both $d$ and $q$ concern this triangulation; neither is invariant under
changes of source.

\subsection{The local product theorem}

Use the global edge-coordinate hyperplanes $\mathcal H$: the equations
$w_e=0$, the three differences of opposite-edge sums in each tetrahedron,
and the twelve face-corner expressions in each tetrahedron.  There are at
most $m+15t$ rows.  Every row is supported on at most four global edges,
has integer coefficients with sum of absolute values at most four, and
comes from one tetrahedron or one edge.  Identifications can cancel or
combine terms but do not increase this bound.

For a set $S$ of edges let $\operatorname{st}(S)$ be the tetrahedra that
meet it, and let $\mathcal H_S$ consist of the rows whose restriction to
$\mathbb R^S$ is nonzero.  Then
\begin{equation}\label{eq:cluster-row-count}
 |\mathcal H_S|\le |S|+15|\operatorname{st}(S)|
                    \le (1+15d)|S|.
\end{equation}

\begin{theorem}[Product of local affine profiles]\label{thm:cluster-product}
Let $F_1,\ldots,F_r$ be the components of $G[F]$.  Fix any admissible
quadrilateral sector $Q$.  In the affine space $b+\mathbb R^F$, the
conditions for membership in its edge-coordinate cone $C_Q$ factor as a
Cartesian product of polyhedra on the separate coordinate sets $F_i$.
The $i$th factor uses only tetrahedra in $\operatorname{st}(F_i)$.

Suppose $w=b+u$ generates a standard vertex normal ray and that
$\operatorname{supp}b\not\subseteq F$.  For each $i$ there are
$s_i=|F_i|$ hyperplanes from $\mathcal H_{F_i}$ such that the system
\begin{equation}\label{eq:cluster-affine}
 R_i u_{F_i}=-a_i,\qquad (a_i)_\ell=\ell(b),
 \quad (R_i)_{\ell,e}=\ell_e,
\end{equation}
is nonsingular and has exactly the target increment as its solution.
Its local reconstruction agrees with $w$ on
$\operatorname{st}(F_i)$ and with the coherent vector elsewhere.
\end{theorem}

\begin{proof}
No tetrahedron can meet two different $F_i$, since its edges would join
those components in $G[F]$.  The edge-coordinate description of normal
surfaces uses only local four-point conditions and face-corner
inequalities: normal matching is automatic when global edge weights
agree.  The constraints in a tetrahedron therefore involve at most one
group of variable increments.  Constraints in an untouched tetrahedron
are constant.  If that constant violates the chosen sector, the slice is
empty; otherwise its nonempty factors have the stated product form.  For
a sector containing the target the latter alternative applies.

Choose an unchanged edge $e$ with $b_e>0$.  The affine slice fixes
$w_e=b_e$, and hence meets the ray $\mathbb R_{>0}w$ at just one point.
Extremality of that ray implies that $w$ is an extreme point of
$C_Q\cap(b+\mathbb R^F)$.  An extreme point of a nonempty Cartesian
product is extreme in every factor: a nontrivial decomposition in one
factor would give such a decomposition of the product point.

At an extreme point of a polyhedron in $\mathbb R^{s_i}$, the active
equalities and inequality boundaries have normals spanning all of
$\mathbb R^{s_i}$.  Otherwise a small displacement in a nonzero common
null direction, of either sign, would stay feasible.  These active
boundaries are all rows from $\mathcal H_{F_i}$.  Select $s_i$
independent ones.  This gives \eqref{eq:cluster-affine}.  Restrictions of
one global row cannot couple different components, which also supplies
a direct block-rank proof.  Finally, each local reconstruction is a
target tetrahedron or a coherent tetrahedron, so it is admissible and
matches globally.  No equality of the target and coherent quadrilateral
types is used.
\end{proof}

\begin{lemma}[The exceptional support is connected]\label{lem:cluster-exception}
If a nonzero admissible standard normal vector spans an extreme ray,
the support of its positive global edge weights is connected in $G$.
Consequently, if $w=b+u$ is a vertex disc with $u_e>0$ exactly on $F$
and $\operatorname{supp}b\subseteq F$, then $G[F]$ is connected.
\end{lemma}

\begin{proof}
If the positive edge support has several components, restrict $w$ to
one component and set every other edge weight to zero.  In each
tetrahedron the resulting weights are either the original weights or
all zero, because a tetrahedron cannot meet two components.  Local
reconstruction gives either the original standard coordinates or zero.
These vectors satisfy matching and the same quadrilateral constraints.
They are nonzero and sum to the original vector, giving a
nonproportional nonnegative decomposition and contradicting extremality.
In the exceptional situation the positive edge support is precisely $F$.
\end{proof}

The lemma concerns the support of $w$, not the support of $w-b$.
Subtracting a coherent vector can disconnect that support.  In
particular it supplies no justification for replacing a component budget
$q(F)\le r$ by the unrestricted assertion $q(F)=1$.

\subsection{A complete algorithm and its cost}

A \emph{profile on $S$} consists of a positive even integral solution of
\eqref{eq:cluster-affine}, the standard coordinates reconstructed on
$\operatorname{st}(S)$, and its change in Euler characteristic.  Every
candidate is checked locally before retention.  Distinct descriptions
that produce the same increments are deduplicated.  The parity
restriction loses no essential discs: adding even increments to the
coherent edge weights preserves even total weight on every face.

Enumerate all connected edge sets $S$ with $1\le |S|\le k$.  For each,
enumerate the $|S|$-subsets of $\mathcal H_S$, solve the square affine
systems, and retain their profiles.  Combine at most $r$ profiles whose
supports are pairwise nonadjacent in $G$ and whose total size is at most
$k$.  Their tetrahedron stars are disjoint, so their local replacement
coordinates assemble without conflict.  The coherent surface is tested
separately.

For an exceptional connected set $S$ containing
$\operatorname{supp}b$, additionally choose $|S|-1$ rows from
$\mathcal H_S$.  Restrict these rows to $\mathbb R^S$, append
$\sum_{e\in S}w_e=1$, and solve whenever the square matrix is
nonsingular.  Keep strictly positive directions and reconstruct a
primitive integral normal vector.  Recompute its actual increments
relative to $b$ before claiming membership in a cancellation class.  The
usual active-ray argument supplies one of these systems for every
exceptional vertex disc.  Lemma~\ref{lem:cluster-exception} ensures that
no disconnected exceptional support is omitted.

\paragraph{Why the positive gate certifies a connected disc.}
For a primitive admissible integral vector $x$, let $J$ be its positive
standard-coordinate support.  If the matching matrix restricted to $J$
has nullity one, each connected component vector is $\lambda x$ for
some $\lambda>0$.  Bezout coefficients for $\gcd(x_j)=1$ show that an
integral vector $\lambda x$ has integral $\lambda$: if
$\sum a_jx_j=1$, then $\lambda=\sum a_j(\lambda x_j)$.  Since the
component coefficients sum to one, there is exactly one component.
Nonzero boundary cohomology modulo two ensures nonempty boundary.  The
classification of compact connected surfaces then makes $\chi(x)=1$
equivalent to a disc.  Its boundary is essential because an essential
simple torus curve has primitive slope, whereas an inessential one has
zero mod-two homology class.  The gate computes the class modulo vertex
coboundaries, and a nonzero supported minor supplies independent exact
rank evidence.  This is the incoming polynomial gate, reused unchanged.

\begin{theorem}[Cancellation-cluster discovery]\label{thm:cluster-complexity}
There is a deterministic exact algorithm which, on a source satisfying
the stated hypotheses, finds an essential normal disc whenever there is
a standard vertex essential disc with
\[
                         |F|\le k,\qquad q(F)\le r.
\]
Writing $B$ for a bound on the binary length of the coherent weights and
source indices, its bit complexity is bounded by
\begin{equation}\label{eq:cluster-cost}
          t^{r+O(1)}(C d^2)^k\operatorname{poly}(k,B)
\end{equation}
for an absolute constant $C$.  A polynomial rank-and-boundary vertex-disc
gate suffices; no bound in the expanded normal-piece count is needed.
The component budget may first be replaced by $\min\{r,k\}$.
For $r=0$ only the coherent case is promised.  The bound may always be
combined, by an ordinary race or by comparing certified input
parameters, with the previous unrestricted cancellation-support bound.
\end{theorem}

\begin{proof}
For a promised nonexceptional disc, Theorem~\ref{thm:cluster-product}
places the true increment on every component among the retained
profiles.  These profiles have pairwise nonadjacent supports, and the
combination stage reconstructs the entire disc.  The exceptional case
was covered above.  Primitive normalization followed by the exact
supported-rank, Euler, and boundary-cohomology gate verifies every
reported positive result.  On a torus boundary a primitive vertex normal
vector of Euler characteristic one and nonzero boundary class modulo two
is an essential disc.  This gate is polynomial in the encoded source and
coordinates.

For counting, a graph of maximum degree $\Delta$ has at most
$m(4\max\{1,\Delta\})^{s-1}$ connected sets of size $s$.  One
elementary enumeration chooses a root, a rooted ordered spanning-tree
shape (at most $4^{s-1}$ choices), and one neighbor for each tree edge;
every connected set has such a description.  With $\Delta\le5d$, this
is at most $m(20d)^{s-1}$.  Equation~\eqref{eq:cluster-row-count} and
$\binom Ns\le(eN/s)^s$ bound the affine systems for one support by
$[e(1+15d)]^s$.  Thus, for some absolute $C_0$, the number of size-$s$
profiles and the work descriptions producing them are at most
$m(C_0d^2)^s$, up to polynomial factors per description.  The
exceptional systems have the same bound: for $s\ge2$ use
$\binom{(1+15d)s}{s-1}\le[2e(1+15d)]^{s-1}$, and handle $s=1$
directly.

For $a\le r$ components of total size $s\le k$, there are
$\binom{s-1}{a-1}$ ordered positive size compositions.  Even before
rejecting incompatible supports, their product lists have length at
most $m^a(C_0d^2)^s$.  Summing over $a,s$ costs a factor at most
$\operatorname{poly}(k)2^k m^r$.  Each exact solve, assembly and
verification costs a fixed polynomial in $t,k,B$ by the bit estimate
below.  Absorb the absolute exponential constants into $C$ to obtain
\eqref{eq:cluster-cost}.  Enumeration, rejected extensions, local
reconstruction, and verification are included in these charges.
\end{proof}

\begin{proposition}[Anchored systems have short binary solutions]
For an affine profile on $s$ edges, the rational increments have a
common denominator of magnitude at most $4^s$ and numerators of
magnitude at most $s4^s\max_e b_e$.  In particular their bit length is
$B+O(s+\log(s+1))$.  The integral profiles retained by the algorithm and
their assembled normal vectors have bit length
$B+O(k+\log(k+1))$ per coordinate.
\end{proposition}

\begin{proof}
Every row of the square matrix $R$ has Euclidean norm at most four.
Hadamard's inequality gives $|\det R|\le4^s$, and every cofactor has
magnitude at most $4^{s-1}$.  The right-hand side has entries of
magnitude at most $4\max b_e$.  The adjugate formula therefore gives
the asserted numerator bound.  Only the single right-hand side carries
the $B$-bit coherent weights; multiplying $B$ by the system order is
unnecessary.  Local reconstruction consists of sums, differences,
minima, and division by two.  The retained profiles have even integral
increments, and disjoint-star assembly takes one profile value per
changed edge and one replacement row per changed tetrahedron.
Exceptional homogeneous rays have integer cofactor representatives of
$O(s)$ bits before reconstruction and primitive normalization, so they
obey the same final bound.
\end{proof}

\subsection{A certified disc with two separated cancellation components}
\label{sec:cluster-counterexample}

Start with the incoming eight-tetrahedron coherent obstruction, attach
four of its prescribed Fibonacci boundary layers, and perform the
$2$--$3$ moves at the currently labelled sites
\[
                        (4,0),\quad(5,0),\quad(12,0).
\]
The resulting solid torus has fifteen tetrahedra, seventeen global edges,
one vertex, maximum tetrahedron-incidence degree $d=8$, and Regina
isomorphism signature
\begin{center}\texttt{pHKzvvPAPQMcbcdlhlhjmmmlnoohqxrorfjidicxak}.\end{center}
The global-edge identifiers used below are not dense indices from one to
seventeen. Its primitive coherent vector has Euler characteristic $-7$.  The
source-bound witness supplied in \texttt{disconnected\_disc.json} is a
primitive standard vertex meridian $D$ with
\[
                 w(D)=b+4\mathbf e_{47}+4\mathbf e_{53}.
\]
The tetrahedron stars are $\operatorname{st}(47)=\{7,13\}$ and
$\operatorname{st}(53)=\{8\}$, so the cancellation support has two
singleton components in $G$.  The witness has $133$ normal pieces and
Euler characteristic one; the coherent vector has $127$ pieces.  All
coordinates, full face pairings, the move history, and an independently
replayable rank-and-boundary certificate are retained.

This is a counterexample to the claim that the cancellation support of
every vertex meridian is connected.  The new algorithm with $k=2,r=1$
exhausts its restricted class without finding a disc, whereas $k=2,r=2$
finds $D$.  The incoming complete weight-space algorithm independently
recovers $D$ from the supplied support $\{47,53\}$.

\begin{proposition}[Exact minimum cancellation support in the retained source]
\label{prop:cluster-fixture-minimum}
For this fifteen-tetrahedron triangulation, the minimum cancellation
support size among standard vertex meridians is
$\kappa_V(\mathcal T)=2$.
\end{proposition}

\begin{proof}[Proof by exact finite algebraic exhaustion]
The independently certified vector $D$ supplies the upper bound.  The
coherent vector has Euler characteristic $-7$, so support size zero is
impossible.  Moreover $b$ is nonzero on fifteen of the seventeen global
edges.  Thus for any singleton $F=\{e\}$ the exceptional case
$\operatorname{supp}b\subseteq F$ is excluded.  By the active-equation
argument in Theorem~\ref{thm:cluster-product}, a vertex meridian with
this support must occur at an offset
\[
       \delta=-\frac{\ell(b)}{\ell_e}>0,
       \qquad \delta\in 2\mathbb Z,
       \qquad \ell\in\mathcal H,\quad \ell_e\ne0.
\]
The script \texttt{audit\_one\_edge\_inventory.py} constructs the dense
global rows directly from the local corner and opposite-sum formulas,
independently of the new profile producer and its sparse hyperplane
builder.  It obtains $119$ distinct rows, $327$ nonzero row--edge
incidences, and $35$ distinct positive even pairs $(e,\delta)$.  Exact
local normal reconstruction rejects $32$ pairs.  The remaining three
are
\[
\begin{array}{c|r|r}
e&\delta&\chi(b+\delta\mathbf 1_e)\\ \hline
42&4&-3\\
47&4&-3\\
53&4&-3
\end{array}
\]
and hence none is a disc.  This exhausts support sizes zero and one,
proving the lower bound.  The source hash, all offsets, all three normal
vectors, and exact Euler characteristics are retained in
\texttt{artifacts/one\_edge\_inventory.json}.
\end{proof}

This finite proof uses the maintained cocycle and coordinate routines
and the incoming local normal decoder; it is independent of the new
candidate enumeration, not an independent reimplementation of every
topological primitive.  It is a reproducible exact computation, not a
proof-assistant formalization or a separate generic negative-certificate
format.  It does not rule out one-component cancellation supports of
arbitrary size, nor does it require the following full vertex census.

For additional finite evidence, Regina 7.4 standard-coordinate tree
traversal and direct double description independently produced the same
$821$ vertex vectors.  Exactly three passed the meridian test; their
$(|F|,q(F))$ pairs are $(2,2),(4,2),(4,2)$.  Every positive meridian
passes the separate exact certificate checker.  The completeness of
this additional census relies on those external enumeration
implementations; the package does not contain an independently checked
negative certificate for absence of all possible one-component vertex
meridians.  The explicit positive counterexample and the restricted
algorithmic theorem do not depend on that stronger census conclusion.

\subsection{Scope of the improvement}

For bounded $d$, fixed $r$, and $k=O(\log t)$,
\eqref{eq:cluster-cost} is polynomial.  For example, bounded $d$,
$r=O(\log t)$, and $k=O((\log t)^2)$ give quasipolynomial discovery.
More generally, when $B$ is polynomial in $t$, the sufficient budget is
\[
       r\log(t+1)+k\log(d+1)=O((\log(t+1))^2).
\]
These implications are conditional on the existence of a suitable
vertex disc in a certified source.  An edge may meet $\Theta(t)$
tetrahedra, $r$ may equal $k$, and no theorem here supplies a source
family with the required bounds.  In the absence of such information,
the unrestricted incoming algorithm can be asymptotically better.

The result is restricted to one global vertex and rank-one cocycles.
With several vertices, the signed cocycle can vary by vertex
coboundaries; independently solving local factors with a fixed vector
$b$ is then incomplete.  The earlier sign-chamber method or a separately
proved treatment of shared cocycle parameters is needed.  Boundary
cohomology in the positive verifier must still be taken modulo all
boundary-vertex coboundaries; a nonzero raw parity vector is insufficient
in that more general setting.
